\chapter*{Conclusion et perspectives}
\addcontentsline{toc}{chapter}{Conclusion et perspectives}
\markboth{Conclusion et perspectives}{Conclusion et perspectives}
\label{sec:conclusion}

Le projet de fin d’études a abouti à la conception et au développement de FESTY, une plateforme numérique dédiée à la gestion événementielle. Cette solution vise à centraliser les principaux services liés à la découverte des événements, à la gestion des partenaires organisateurs, à la billetterie, au paiement, à la revente des tickets, au contrôle d’accès et à la supervision administrative.

Tout au long de ce projet, nous avons suivi une démarche progressive basée sur la méthodologie Agile Scrum. Cette organisation nous a permis de découper le travail en quatre sprints cohérents, de prioriser les besoins et de construire progressivement les différentes fonctionnalités de la plateforme.

Dans un premier temps, nous avons présenté le contexte général du projet, l’organisme d’accueil TEKSIGHT, la problématique étudiée ainsi que l’étude de l’existant. Cette analyse a mis en évidence la fragmentation des solutions actuelles et le besoin d’une plateforme centralisée capable de réunir la découverte des événements, la gestion des billets, le paiement, la revente, le contrôle d’accès et le suivi administratif.

Nous avons ensuite proposé une solution composée d’une application mobile, d’un espace partenaire et d’un back-office d’administration. Cette solution permet aux utilisateurs de découvrir les événements, de personnaliser leur expérience, de réserver et payer leurs billets, de consulter leurs tickets et d’utiliser la marketplace de revente. Elle permet également aux partenaires de gérer leurs événements, leur line-up artistique, leurs sessions de scan et leur activité. Enfin, elle offre à l’administrateur des outils de contrôle de l’offre événementielle, de gestion du catalogue artistique et de suivi des opérations sensibles.

Sur le plan fonctionnel, le premier sprint a permis de mettre en place l’inscription, l’authentification, la vérification du numéro de téléphone, la gestion du profil utilisateur et l’accès sécurisé au back-office administrateur. Le deuxième sprint a été consacré à la gestion des partenaires et des événements. Le troisième sprint a introduit l’exploration, la personnalisation, les recommandations, les favoris, le suivi des stars, le chat événementiel et la gestion du catalogue artistique. Enfin, le quatrième sprint a complété le parcours utilisateur avec la billetterie, le paiement, la gestion des tickets, le remboursement, la marketplace de revente, le scan des billets et la supervision financière.

À l’issue de ce travail, nous avons réussi à concevoir une plateforme cohérente, évolutive et adaptée aux besoins des différents acteurs du domaine événementiel. FESTY apporte une réponse aux problèmes de dispersion des services en regroupant dans un même environnement les fonctionnalités nécessaires aux participants, aux organisateurs, aux agents de scan et aux administrateurs. Cette solution contribue ainsi à améliorer l’expérience utilisateur, à simplifier la gestion des événements et à renforcer la sécurité des opérations liées aux billets et aux paiements.

En termes de perspectives, plusieurs améliorations peuvent être envisagées afin d’enrichir davantage la plateforme. Nous prévoyons notamment d’améliorer le moteur de recommandation afin de proposer des événements plus pertinents selon les préférences, les interactions, les stars suivies et la localisation des utilisateurs. À moyen terme, l’intégration de techniques d’intelligence artificielle pourrait permettre d’affiner davantage la personnalisation.

Une autre perspective consiste à renforcer les notifications en temps réel afin d’informer rapidement les utilisateurs des changements liés aux événements, des confirmations de paiement, des remboursements, des mises à jour de tickets et des messages dans le chat.

Nous envisageons également d’enrichir le tableau de bord partenaire avec des indicateurs plus détaillés sur les ventes, le taux de remplissage, les revenus, les tickets scannés et la performance des événements. Ces indicateurs aideraient les organisateurs à mieux suivre leur activité et à prendre des décisions plus efficaces.

La supervision administrative pourrait aussi être améliorée par l’ajout de mécanismes avancés de détection des opérations suspectes, notamment au niveau des paiements, des remboursements, de la revente des tickets et des accès scannés. Enfin, l’ergonomie des interfaces et la scalabilité de l’architecture pourraient être renforcées afin d’accompagner l’évolution future de FESTY.